\documentclass[11pt,a4paper]{article}
\usepackage[margin=2.5cm]{geometry}
\usepackage{amsmath,amssymb,amsthm}
\usepackage{algorithm,algorithmic}
\usepackage{bm}
\usepackage{hyperref}
\usepackage{enumitem}

\newcommand{\R}{\mathbb{R}}
\newcommand{\E}{\mathbb{E}}
\newcommand{\Var}{\operatorname{Var}}
\newcommand{\sign}{\operatorname{sign}}
\newcommand{\zscore}{\operatorname{zscore}}
\newcommand{\med}{\operatorname{med}}
\DeclareMathOperator*{\argmin}{argmin}

\theoremstyle{definition}
\newtheorem{definition}{Definition}
\newtheorem{remark}{Remark}

\title{GeoXGB: Mathematical Architecture of\\Geometry-Aware Gradient Boosting}
\author{}
\date{}

\begin{document}
\maketitle

\section{Setup and Notation}

Let $X \in \R^{n \times d}$ be the training feature matrix and $y \in \R^n$ the
target vector.  Throughout, subscript $i$ indexes samples and superscript $(j)$
indexes features.  We write $x_i \in \R^d$ for the $i$-th row of $X$.

\subsection{Whitening}

HVRT operates on \emph{whitened} (z-scored) features.  The whitener computes
per-column statistics:
\begin{equation}
    z_i^{(j)} = \frac{x_i^{(j)} - \mu^{(j)}}{\sigma^{(j)} + \epsilon},
    \qquad
    \mu^{(j)} = \frac{1}{n}\sum_{i=1}^{n} x_i^{(j)},
    \quad
    \sigma^{(j)} = \sqrt{\frac{1}{n}\sum_{i=1}^{n} \bigl(x_i^{(j)} - \mu^{(j)}\bigr)^2},
\end{equation}
with $\epsilon = 10^{-8}$.  This gives the whitened matrix $Z \in \R^{n \times d}$
with each column approximately $\mathcal{N}(0,1)$.  The whitener parameters
$(\mu, \sigma)$ are computed \emph{once} at the initial HVRT fit and frozen
for all subsequent refits.

\subsection{Inverse Transform}

Synthetic samples generated in $z$-space are mapped back to the original feature
space via
\begin{equation}
    \hat{x}^{(j)} = z^{(j)} \cdot \sigma^{(j)} + \mu^{(j)}.
\end{equation}


\section{The Geometry Target}

\subsection{Pairwise Cooperation Target (HVRT)}

For the default HVRT partitioner ($d \le 50$), the geometry target $T \in \R^n$
is the \emph{pairwise cooperation statistic}:
\begin{equation}
    T_i = \zscore\!\Biggl(\,\sum_{a < b} \zscore\bigl(z_i^{(a)} \cdot z_i^{(b)}\bigr)\Biggr),
\end{equation}
where $\zscore(v) = (v - \bar{v}) / \sigma_v$ standardises a vector to zero mean
and unit variance (population).

\begin{remark}[Algebraic identity]
$T$ is proportional to the second elementary symmetric polynomial:
\[
    \sum_{a < b} z_i^{(a)} z_i^{(b)}
    = \frac{1}{2}\Bigl[\Bigl(\sum_j z_i^{(j)}\Bigr)^{\!2} - \sum_j \bigl(z_i^{(j)}\bigr)^2\Bigr]
    = \frac{S_i^2 - Q_i}{2}
    = e_2(\bm{z}_i),
\]
where $S_i = \sum_j z_i^{(j)}$ and $Q_i = \|\bm{z}_i\|_2^2$.  The z-score
normalisation within each pair column absorbs the scale but preserves the
rank ordering, so $T$ partitions samples along iso-$e_2$ hyperboloids
in $z$-space.  Computational cost: $\mathcal{O}(n \cdot d^2)$.
\end{remark}

\begin{remark}[Noise invariance]
If each feature is corrupted by independent additive noise
$\tilde{z}^{(j)} = z^{(j)} + \varepsilon^{(j)}$ with
$\E[\varepsilon^{(j)}] = 0$, then
\[
    \E[\tilde{S}^2 - \tilde{Q}]
    = \E[S^2 - Q] + d\sigma_\varepsilon^2 - d\sigma_\varepsilon^2
    = \E[S^2 - Q].
\]
The additive noise bias cancels exactly because both $S^2$ and $Q$ are
degree-2 in the features.  This is the HVRT noise invariance theorem.
\end{remark}

\subsection{Alternative Targets}

\begin{itemize}[nosep]
    \item \textbf{HART}: $T_i = \zscore\!\bigl(\frac{1}{2}(\|z_i\|_1^2 - \|z_i\|_2^2)\bigr) = \zscore\!\bigl(\sum_{a<b}|z_i^{(a)}||z_i^{(b)}|\bigr)$.
        Degree-2, noise-invariant, $\mathcal{O}(n \cdot d)$.
    \item \textbf{PyramidHART}: $T_i = \zscore(|S_i| - \|z_i\|_1)$.
        Degree-1, outlier-cancelling, $\mathcal{O}(n \cdot d)$.
    \item \textbf{FastHART} ($d > 50$): $T_i = \zscore(S_i)$.
        Degree-1, $\mathcal{O}(n \cdot d)$.
\end{itemize}


\section{Blended Target and Partition Tree}

\subsection{Y-Weight Blending}

The pure geometry target $T$ is blended with the $y$-signal (raw targets on
initial fit, or residual gradients on refit) to bias the partition tree towards
regions that are both geometrically interesting \emph{and} predictively
important.

\begin{align}
    y_{\text{norm}} &= \frac{y - y_{\min}}{y_{\max} - y_{\min}} \in [0,1]^n, \\
    y_{\text{dev}}  &= \bigl|y_{\text{norm}} - \med(y_{\text{norm}})\bigr|, \\
    y_{\text{comp}} &= \zscore(y_{\text{dev}}), \\
    t_{\text{blend}} &= \zscore\!\bigl((1 - w_y)\,\zscore(T) + w_y\,y_{\text{comp}}\bigr),
\end{align}
where $w_y \in [0,1]$ is the \texttt{y\_weight} parameter.

\begin{remark}
$y_{\text{comp}}$ measures distance from the median response: extreme $y$
values (both high and low) get large $y_{\text{comp}}$.  This is deliberate ---
the partition tree should separate extreme-response regions regardless of sign,
because both tails carry the most predictive information for gradient boosting.
\end{remark}

\subsection{Partition Tree}

The blended target $t_{\text{blend}}$ is used as the splitting criterion for a
variance-reduction decision tree on $Z$.  The tree greedily minimises
\begin{equation}
    \argmin_{j, \theta} \;\;
    \sum_{i \in L} (t_i - \bar{t}_L)^2 + \sum_{i \in R} (t_i - \bar{t}_R)^2,
    \qquad
    L = \{i : z_i^{(j)} \le \theta\},\;\; R = \{i : z_i^{(j)} > \theta\},
\end{equation}
producing a set of $P$ leaf partitions $\{\mathcal{P}_1, \dots, \mathcal{P}_P\}$.
Each sample $i$ is assigned to exactly one partition $p(i)$.

The partition tree operates on binned features (16 bins) for speed, using
histogram-based split enumeration.  The tree depth and leaf sizes are
auto-tuned: $\texttt{min\_samples\_leaf} = 10$,
$\texttt{max\_partitions} = \min(1500,\; 3n / 20)$.


\section{Reduce--Expand Resampling}

Given the partition structure $\{(\mathcal{P}_p, Z_p)\}_{p=1}^{P}$, HVRT
performs a \emph{reduce} step (discard uninformative samples) followed by an
optional \emph{expand} step (generate synthetic samples).

\subsection{Budget Allocation}

Each partition $p$ receives a budget $b_p$ of how many samples to retain
(or generate).  Two modes:
\begin{itemize}[nosep]
    \item \textbf{Proportional} (default): $b_p \propto |\mathcal{P}_p|$.
    \item \textbf{Variance-weighted}: $b_p \propto \bar{v}_p$ where
        $\bar{v}_p = \frac{1}{|\mathcal{P}_p|}\sum_{i \in \mathcal{P}_p} \frac{1}{d}\sum_j |z_i^{(j)}|$
        is the mean absolute deviation within partition $p$.
\end{itemize}
The total budget is $n_{\text{keep}} = \lfloor n \cdot r_{\text{eff}} \rfloor$,
where $r_{\text{eff}}$ is the effective reduce ratio (see \S\ref{sec:noise}).

\subsection{Variance-Ordered Reduction}

Within each partition $p$ of size $n_p$ with budget $b_p < n_p$, the
\textbf{variance-ordered} selector retains samples that maximise geometric
diversity:

\begin{enumerate}[nosep]
    \item Compute the $n_p \times n_p$ pairwise squared-distance matrix
        $D_{ij} = \|z_i - z_j\|_2^2$ for $i, j \in \mathcal{P}_p$.
    \item For each $i$, find the $k$-nearest neighbours ($k = 5$) and compute
        $v_i = \Var(\{D_{ij}\}_{j \in \text{kNN}(i)})$ --- the local distance variance.
    \item Sort by $v_i$ descending and retain the top $b_p$ samples.
\end{enumerate}

High local-distance-variance points lie near partition boundaries or in
heterogeneous regions of $z$-space.  Retaining them preserves the geometric
``skeleton'' of each partition while discarding interior samples that are
well-represented by their neighbours.

\begin{remark}[Scalability]
When $n_p > 400$, a random subsample of $400$ candidates is drawn first,
variance-ordering is applied to the subsample, and the selected indices are
mapped back.  When $b_p / n_p \ge 0.65$, the selection degenerates to random
sampling (Fisher--Yates) since the quality gain over random is negligible at
high retention fractions.
\end{remark}

\subsection{Expansion: Synthetic Sample Generation}

When $\texttt{expand\_ratio} > 0$ or $n < \texttt{min\_train\_samples}$,
HVRT generates $n_{\text{syn}} = \lfloor n \cdot e \cdot \hat{\nu} \rfloor$
synthetic samples in $z$-space, where $e$ is the expand ratio and
$\hat{\nu}$ is the noise modulation (\S\ref{sec:noise}).

Synthetic samples are generated per-partition via kernel density estimation.
For the \textbf{Epanechnikov} kernel (default): within partition $p$ with
samples $\{z_i\}_{i \in \mathcal{P}_p}$, draw a centre $z_c$ uniformly at
random from the partition's training points, then sample
\begin{equation}
    z_{\text{syn}} = z_c + h_p \cdot u, \qquad u \sim \text{Epanechnikov}(d),
\end{equation}
where $h_p$ is the per-partition bandwidth (Scott's rule:
$h_p = n_p^{-1/(d+4)}\,\hat{\sigma}_p$).

The synthetic $z$-vectors are inverse-transformed back to feature space,
giving $X_{\text{syn}} \in \R^{n_{\text{syn}} \times d}$.

\subsection{Y-Assignment for Synthetic Samples}
\label{sec:knn}

Synthetic samples have features but no target values.  These are assigned
via \textbf{inverse-distance-weighted $k$-NN} ($k = 3$) in $z$-space:
\begin{equation}
    \hat{y}_{\text{syn}} = \frac{\sum_{j \in \text{kNN}} w_j \, y_j}{\sum_{j \in \text{kNN}} w_j},
    \qquad
    w_j = \frac{1}{\|z_{\text{syn}} - z_j\|_2 + \epsilon},
\end{equation}
where the nearest-neighbour search is over the \emph{reduced} real training
set $\{(z_i, y_i)\}_{i \in \mathcal{R}}$ and $\epsilon = 10^{-8}$.

The final training matrix for each boosting interval is:
\begin{equation}
    \tilde{X} = \begin{bmatrix} X_{\text{red}} \\ X_{\text{syn}} \end{bmatrix}
    \in \R^{(n_{\text{red}} + n_{\text{syn}}) \times d},
    \qquad
    \tilde{y} = \begin{bmatrix} y_{\text{red}} \\ \hat{y}_{\text{syn}} \end{bmatrix}.
\end{equation}


\section{The Boosting Loop}

\subsection{Initialisation}

\begin{enumerate}[nosep]
    \item Compute initial prediction:
        \begin{itemize}[nosep]
            \item \textbf{Regression}: $\hat{y}_0 = \bar{y}$ (mean of training targets).
            \item \textbf{Classification}: $\hat{y}_0 = \log\frac{\bar{y}}{1 - \bar{y}}$ (log-odds).
        \end{itemize}
    \item Set all predictions: $\hat{y}_i^{(0)} = \hat{y}_0$ for all $i$.
    \item Run initial HVRT fit on $(X, y)$ using raw targets as the $y$-signal.
    \item Reduce and optionally expand, producing $(\tilde{X}^{(0)}, \tilde{y}^{(0)})$.
    \item Set resampled predictions:
        $\tilde{p}_i^{(0)} = \hat{y}_0$ for all $i \in \tilde{X}^{(0)}$.
\end{enumerate}

\subsection{Round-by-Round Gradient Descent}

For round $t = 1, \dots, T$ (\texttt{n\_rounds}):

\paragraph{Step 1: Compute gradients on the resampled set.}
\begin{align}
    g_i^{(t)} &= \tilde{y}_i^{(t-1)} - \tilde{p}_i^{(t-1)}
    && \text{(regression, squared error)} \\
    g_i^{(t)} &= \tilde{y}_i^{(t-1)} - \sigma(\tilde{p}_i^{(t-1)})
    && \text{(classification, log-loss)}
\end{align}
where $\sigma(v) = 1/(1 + e^{-v})$ is the sigmoid function.

\paragraph{Step 2: Fit a weak learner tree on the resampled data.}

A decision tree $h_t$ is built on $(\tilde{X}^{(t-1)}, g^{(t)})$ using:
\begin{itemize}[nosep]
    \item Pre-computed bin edges from the \emph{full} dataset (64 bins).
    \item Random split strategy (ExtraTrees-style): for each candidate feature,
        the split threshold is drawn uniformly from the bin edges.
    \item Variance-reduction splitting on the gradient values.
    \item \texttt{max\_depth}, \texttt{min\_samples\_leaf} control tree complexity.
\end{itemize}

\paragraph{Step 3: Update predictions.}
\begin{align}
    \tilde{p}_i^{(t)} &= \tilde{p}_i^{(t-1)} + \eta \cdot h_t(\tilde{x}_i)
    && \text{(resampled set)} \\
    \hat{y}_i^{(t)} &= \hat{y}_i^{(t-1)} + \eta \cdot h_t(x_i)
    && \text{(full training set)}
\end{align}
where $\eta$ is the \texttt{learning\_rate}.

\begin{remark}[Dual prediction tracking]
GeoXGB maintains \emph{two} prediction vectors:
$\tilde{p}$ on the resampled set (used for gradient computation on the
reduced+synthetic data) and $\hat{y}$ on the full training set (used for
gradient computation at refit boundaries and for the final model output).
Each tree is evaluated on both sets every round: $h_t(\tilde{X})$ and $h_t(X)$.
This is $\mathcal{O}(n)$ per round, not $\mathcal{O}(n_{\text{red}})$.
\end{remark}


\section{Refits: Geometry Re-Survey}
\label{sec:refit}

Every \texttt{refit\_interval} rounds (e.g.\ every 50 rounds), the partition
geometry is re-surveyed using the \emph{current gradient landscape}.  This is
the critical mechanism that couples the static geometry to the evolving
loss surface.

\subsection{Refit Trigger}

At round $t$ where $t \bmod R = 0$ (with $R = \texttt{refit\_interval}$):

\begin{enumerate}[nosep]
    \item Compute full-data gradients:
        $g_i^{(\text{full})} = y_i - \hat{y}_i^{(t)}$ (regression) or
        $g_i^{(\text{full})} = y_i - \sigma(\hat{y}_i^{(t)})$ (classification).
    \item Pass $g^{(\text{full})}$ as the new $y$-signal to HVRT refit.
\end{enumerate}

\subsection{What Changes at Refit}

The HVRT refit is a \emph{fast path} that reuses:
\begin{itemize}[nosep]
    \item The whitened matrix $Z$ (computed once at initial fit).
    \item The binned feature matrix (computed once at initial fit).
    \item The geometry target $T$ (computed once at initial fit).
\end{itemize}

What changes:
\begin{enumerate}[nosep]
    \item The $y$-signal in the blend is now the gradient vector
        $g^{(\text{full})}$ instead of the raw targets $y$.
    \item The blended target becomes:
        \begin{equation}
            t_{\text{blend}}^{(t)} = \zscore\!\bigl(
                (1 - w_y^{(t)})\,T + w_y^{(t)}\,y_{\text{comp}}(g^{(\text{full})})
            \bigr),
        \end{equation}
        where $y_{\text{comp}}(g) = \zscore(|g_{\text{norm}} - \med(g_{\text{norm}})|)$
        measures gradient extremality.
    \item The partition tree is rebuilt on $Z$ with the new blended target.
    \item The expander (KDE parameters) is re-fitted if partitions changed.
    \item Reduction and expansion produce a new $(\tilde{X}^{(t)}, \tilde{y}^{(t)})$.
\end{enumerate}

\subsection{Geometric Interpretation of Refit}

\begin{definition}[Refit as gradient-geometry alignment]
At round 0, the partition tree separates $z$-space along iso-$T$ level sets
blended with raw-target extremality.  The tree knows which samples have high
pairwise feature cooperation but knows nothing about which regions the
boosted ensemble struggles with.

At refit round $t$, the gradient $g^{(\text{full})}$ encodes exactly where
the ensemble is wrong: large $|g_i|$ means sample $i$ is poorly predicted.
The blended target $t_{\text{blend}}^{(t)}$ now separates $z$-space into
regions that are both geometrically coherent \emph{and} where the ensemble
has high residual error.

This means:
\begin{itemize}[nosep]
    \item Partitions \emph{drift} toward hard-to-predict regions of $z$-space.
    \item The reduced set preferentially retains samples with high residuals.
    \item Synthetic samples are generated near geometrically interesting points
        that the ensemble cannot yet fit.
    \item As the ensemble improves, gradients shrink, $y_{\text{comp}}$ becomes
        more uniform, and the partition tree reverts toward pure geometry.
\end{itemize}
\end{definition}

\subsection{Adaptive Y-Weight}
\label{sec:adaptive_yw}

When \texttt{adaptive\_y\_weight} is enabled, the effective $y$-weight
is scaled by the alignment between the geometry target and the current
residuals:
\begin{equation}
    w_y^{(t)} = w_y \cdot |\rho(T, g^{(\text{full})})|,
\end{equation}
where $\rho$ is the Pearson correlation.  When the gradients are orthogonal
to the geometry ($\rho \approx 0$), $w_y^{(t)} \to 0$ and the partition
tree becomes pure geometry.  When the gradients are strongly aligned with
the geometry ($|\rho| \approx 1$), the full $y$-weight is applied.

\begin{remark}
This prevents the gradient signal from corrupting the geometry when the
residuals are pure noise (late in training when the ensemble has converged).
It is the adaptive mechanism that allows GeoXGB to self-regulate the
geometry--gradient coupling.
\end{remark}


\section{Noise Modulation}
\label{sec:noise}

\subsection{Signal-to-Noise Estimator}

At each resample event, GeoXGB estimates the local signal-to-noise ratio
$\hat{\nu} \in [0, 1]$:

\begin{enumerate}[nosep]
    \item Subsample $m = \min(500, n)$ probe points.
    \item For each probe $i$, find its $k$-nearest neighbours in $z$-space
        ($k = \min(10, n/30)$) and compute the local mean
        $\bar{y}_i^{(\text{local})} = \frac{1}{k}\sum_{j \in \text{kNN}(i)} y_j$.
    \item Compute local-mean variance $\sigma_{\text{local}}^2 = \Var(\bar{y}^{(\text{local})})$
        and total $y$-variance $\sigma_y^2 = \Var(y)$.
    \item SNR ratio: $\text{snr} = \sigma_{\text{local}}^2 / \sigma_y^2$.
    \item $\hat{\nu} = \text{clamp}\!\bigl(\frac{\text{snr} - 0.15}{0.45},\; 0,\; 1\bigr)$.
\end{enumerate}

\begin{remark}
$\hat{\nu} = 1$ means pure signal (local averages explain all $y$-variance).
$\hat{\nu} = 0$ means pure noise (local averages are constant --- $y$ has
no spatial structure in $z$-space).
\end{remark}

\subsection{How Noise Modulation Affects the Loop}

\begin{itemize}[nosep]
    \item \textbf{Reduce ratio}: $r_{\text{eff}} = r + (1 - \hat{\nu})(1 - r)$.
        Noisier data $\to$ keep more samples (less aggressive reduction).
    \item \textbf{Expand count}: $n_{\text{syn}} = \lfloor n \cdot e \cdot \hat{\nu} \rfloor$.
        Noisier data $\to$ generate fewer synthetic samples.
    \item \textbf{Noise guard}: when $\hat{\nu} < 0.05$ and $n < \texttt{min\_train\_samples}$,
        skip the refit entirely (the gradient is too noisy to guide the geometry).
\end{itemize}


\section{Block Cycling (Large-$n$ Scalability)}

When $n > n_{\text{block}}$ (auto-determined or user-specified), the full
training set is divided into non-overlapping blocks of size $n_{\text{block}}$.

\begin{itemize}[nosep]
    \item A deterministic LCG-based permutation shuffles sample indices.
    \item At each refit boundary, the boosting loop advances to the next block.
    \item When all blocks are exhausted, a new epoch begins with a fresh
        permutation.
    \item Each block gets a \emph{fresh} HVRT fit (whitening, geometry target,
        partition tree) since different blocks cover different geometric regions.
\end{itemize}

\begin{equation}
    n_{\text{block}} = \max\!\bigl(2000,\; \lfloor\sqrt{n} \cdot 15 \cdot \text{ri\_scale}\rfloor\bigr),
    \qquad
    \text{ri\_scale} = \text{clamp}\!\Bigl(\frac{R}{50},\; 0.5,\; 2.0\Bigr).
\end{equation}

\begin{remark}
Block cycling makes all per-round costs $\mathcal{O}(n_{\text{block}})$
instead of $\mathcal{O}(n)$.  The $\sqrt{n}$ scaling ensures blocks grow
sub-linearly: at $n = 100{,}000$ with $R = 50$, $n_{\text{block}} \approx 4{,}743$.
\end{remark}


\section{Complete Algorithm}

\begin{algorithm}[H]
\caption{GeoXGB Fit (Regression, Squared Error)}
\begin{algorithmic}[1]
\REQUIRE $X \in \R^{n \times d}$, $y \in \R^n$, hyperparameters $\eta, R, r, e, w_y, T_{\max}$
\STATE $\hat{y}_0 \gets \bar{y}$; \quad $\hat{\bm{y}} \gets \hat{y}_0 \cdot \bm{1}_n$
\STATE Compute whitener: $Z \gets \text{whiten}(X)$
\STATE Compute geometry target: $T \gets \text{pairwise\_target}(Z)$
\STATE $t_{\text{blend}} \gets \zscore\bigl((1-w_y) T + w_y\, y_{\text{comp}}(y)\bigr)$
\STATE Build partition tree on $(Z, t_{\text{blend}})$ $\to$ partitions $\{\mathcal{P}_p\}$
\STATE Fit KDE expander per partition
\STATE Reduce: $\mathcal{R} \gets \text{variance\_ordered}(\{\mathcal{P}_p\}, \lfloor n r \rfloor)$
\STATE Expand: $\mathcal{S} \gets \text{KDE\_sample}(\lfloor n e \rfloor)$; \quad $\hat{y}_{\text{syn}} \gets \text{kNN}_{k=3}(Z_\mathcal{S}, Z_\mathcal{R}, y_\mathcal{R})$
\STATE $\tilde{X} \gets [X_\mathcal{R}; X_\mathcal{S}]$; \quad $\tilde{y} \gets [y_\mathcal{R}; \hat{y}_{\text{syn}}]$
\STATE $\tilde{\bm{p}} \gets \hat{y}_0 \cdot \bm{1}_{|\tilde{X}|}$
\FOR{$t = 1$ \textbf{to} $T_{\max}$}
    \IF{$t > 0$ and $t \bmod R = 0$}
        \STATE $g^{(\text{full})} \gets y - \hat{\bm{y}}$ \hfill\COMMENT{Full-data gradients}
        \STATE $w_y^{(t)} \gets w_y \cdot |\rho(T, g^{(\text{full})})|$ \hfill\COMMENT{Adaptive y-weight}
        \STATE $t_{\text{blend}}^{(t)} \gets \zscore\bigl((1-w_y^{(t)}) T + w_y^{(t)}\, y_{\text{comp}}(g^{(\text{full})})\bigr)$
        \STATE Rebuild partition tree on $(Z, t_{\text{blend}}^{(t)})$ \hfill\COMMENT{Refit: same $Z$, $T$; new blend}
        \STATE Re-fit KDE if partitions changed
        \STATE Re-reduce, re-expand $\to$ new $(\tilde{X}, \tilde{y})$
        \STATE Sync $\tilde{\bm{p}}$ from $\hat{\bm{y}}$ for real samples, $\text{predict}(\tilde{X}_{\text{syn}})$ for synthetic
        \STATE $\tilde{y}_{\text{real}} \gets \tilde{\bm{p}}_{\text{real}} + g^{(\text{full})}_\mathcal{R}$ \hfill\COMMENT{Reconstruct targets}
    \ENDIF
    \STATE $g \gets \tilde{y} - \tilde{\bm{p}}$ \hfill\COMMENT{Gradients on resampled set}
    \STATE Fit tree $h_t$ on $(\tilde{X}, g)$ with random splits, depth=$D$
    \STATE $\tilde{\bm{p}} \gets \tilde{\bm{p}} + \eta \cdot h_t(\tilde{X})$
    \STATE $\hat{\bm{y}} \gets \hat{\bm{y}} + \eta \cdot h_t(X)$
\ENDFOR
\STATE $F(x) = \hat{y}_0 + \eta \sum_{t=1}^{T_{\max}} h_t(x)$ \hfill\COMMENT{Ensemble prediction}
\end{algorithmic}
\end{algorithm}


\section{Geometric Evolution Across $n$}

\subsection{Small $n$ ($n \lesssim 500$)}

\begin{itemize}[nosep]
    \item Few partitions (10--50).  Each partition covers a broad region of $z$-space.
    \item Variance-ordered reduction retains most samples (at $r = 0.7$, keeping 350 of 500).
    \item Synthetic expansion fills gaps: with $e = 0.1$, generates $\sim$50 synthetic points.
    \item The kNN $y$-assignment uses the full reduced set as reference --- accurate because
        partition interiors are well-sampled.
    \item Refits shift partitions significantly: $g^{(\text{full})}$ on 500 points is noisy,
        so the blended target changes substantially each refit.
    \item Net effect: the geometry acts as a \emph{regulariser} --- the partition structure
        prevents the boosted trees from overfitting to individual noisy samples.
\end{itemize}

\subsection{Moderate $n$ ($n \sim 1{,}000\text{--}10{,}000$)}

\begin{itemize}[nosep]
    \item More partitions (50--500).  Finer geometry captures local feature interactions.
    \item Reduction is meaningful: at $r = 0.7$, discarding 30\% of 5{,}000 = 1{,}500
        uninformative interior points.  The reduced set of 3{,}500 preserves the geometric
        skeleton.
    \item Synthetic samples concentrate where the KDE has high density --- near partition
        centroids.  These provide ``consensus'' targets that smooth the gradient signal.
    \item Refits are stable: $g^{(\text{full})}$ on 5{,}000+ points has low variance,
        so partition drift is gradual and predictable.
    \item This is the regime where GeoXGB outperforms XGBoost most consistently:
        the geometry captures feature-interaction structure that pure tree splitting
        must discover through many sequential splits.
\end{itemize}

\subsection{Large $n$ ($n \gtrsim 20{,}000$)}

\begin{itemize}[nosep]
    \item Block cycling activates: each refit sees a different $n_{\text{block}}$-sized
        slice of the data.
    \item Each block gets a \emph{fresh} whitener and geometry target.  The partitions
        are block-local, not global.
    \item The ensemble sees different geometric views at each refit, providing
        implicit regularisation through geometric diversity.
    \item At $n = 100{,}000$, $n_{\text{block}} \approx 4{,}700$.  Over $T = 1{,}000$ rounds
        with $R = 50$, the ensemble trains on $\sim$20 blocks $\times$ 3{,}300 reduced
        samples = 66{,}000 unique samples --- progressive coverage of the full dataset.
    \item Prediction cost is $\mathcal{O}(n_{\text{block}})$ per round instead of
        $\mathcal{O}(n)$, making the total fit cost $\mathcal{O}(T \cdot n_{\text{block}} \cdot d)$.
\end{itemize}

\subsection{Summary: What Each Component Does}

\begin{center}
\begin{tabular}{lll}
\hline
\textbf{Component} & \textbf{Mathematical Role} & \textbf{Effect on Learning} \\
\hline
Whitening & $z = (x - \mu)/\sigma$ & Isotropic feature space \\
Geometry target $T$ & $e_2(z)$ level sets & Feature-interaction structure \\
Y-weight blend & $(1-w)T + w\,|y - \med(y)|$ & Gradient--geometry coupling \\
Partition tree & Variance-reduction splits on $t_{\text{blend}}$ & Spatial decomposition \\
Reduction & Retain high-variance boundary points & Remove redundant interior \\
Expansion & Per-partition KDE sampling & Smooth gradient signal \\
kNN $y$-assign & IDW interpolation in $z$-space & Consistent synthetic targets \\
Refit & New blend with current gradients & Track evolving loss surface \\
Adaptive $w_y$ & $w_y \cdot |\rho(T, g)|$ & Self-regulate coupling \\
Noise modulation & Local SNR estimate & Gate expansion/reduction \\
Block cycling & Rotate $n_{\text{block}}$-sized slices & Scalable geometric diversity \\
\hline
\end{tabular}
\end{center}


\end{document}
